This section constructs the integral word order of Theorem~\ref{thm:hecke-word-descent} and an explicit type of denominator certificate. The Hecke normalization is that of \citet[Section 3.1, pp.~20--21]{BSM2013}.

\begin{proof}[Proof of Theorem~\ref{thm:hecke-word-descent}]
The four substitutions $v\mapsto v,-v,v^{-1},-v^{-1}$ are distinct automorphisms of $K_v$ fixing $D$. Conversely, $v$ satisfies
\[
 v^4-(D-2)v^2+1=0.
\]
It follows that $K_v/F_D$ is Galois of degree four, with precisely these automorphisms. Also
\[
 v^{-1}=(D-2)v-v^3.
\]
The quartic relation and this identity show that $1,v,v^2,v^3$ span $R_v$ over $S_D$. Their independence follows from $[K_v:F_D]=4$. Hence $R_v$ is free of rank four over $S_D$, and is flat. Moreover $R_v\otimes_{S_D}F_D=K_v$.

There are two commuting semilinear involutions of the Hecke algebra over $K_v$:
\[
 \alpha(v)=-v,\quad \alpha(T_s)=-T_s;
 \qquad
 \beta(v)=v^{-1},\quad \beta(T_s)=T_s^{-1}.
\]
They preserve the quadratic and braid relations. Indeed the quadratic relation gives
\[
 T_s^{-1}=T_s-v+v^{-1},
\]
and inversion of the braid relations gives the corresponding relations for the inverses. The two involutions commute on scalars and on every $T_s$, where both compositions give $-T_s^{-1}$. Apply each involution diagonally to the two tensor factors. The map $\alpha$ negates both $T_s+v^{-1}$ and $T_s-v$, whereas $\beta$ fixes both. Therefore both involutions fix $P_s$, and hence fix every word.

Choose a $K_v$-basis $w_1,\ldots,w_d$ of $A_{K_v}$ from its spanning set of words. For any word $w$, write
\[
 w=\sum_{i=1}^d a_iw_i,\qquad a_i\in K_v.
\]
Applying $\alpha$ and $\beta$ and using uniqueness of these coordinates shows that each $a_i$ is fixed by the Galois group. Thus $a_i\in F_D$. The same words are consequently an $F_D$-basis of $A_{F_D}$, and scalar extension gives $A_{K_v}$. This proves the field isomorphism and the assertion about multiplication coefficients, because a product of two words is again a word.

The standard Hecke basis makes $M_{R_v}$ free of rank $(r!)^2$ over $R_v$, and therefore free of rank $4(r!)^2$ over $S_D$. The module $A_{S_D}$ is an $S_D$-submodule of this finite free module. Since $S_D$ is Noetherian, it is finite, and its embedding proves torsion-freeness. Word concatenation makes it a unital algebra.

The word basis identifies $A_{F_D}$ with $F_D^d$. Tensor the inclusion $A_{S_D}\subseteq A_{F_D}$ with the flat $S_D$-module $R_v$. This gives an injection
\[
 R_v\otimes_{S_D}A_{S_D}
 \longrightarrow R_v\otimes_{S_D}A_{F_D}
 \cong K_v\otimes_{F_D}A_{F_D}
 \cong A_{K_v}.
\]
Its image is exactly the $R_v$-span of the words, namely $A_{R_v}$. This proves the integral isomorphism. It uses flatness and the field calculation; no unramified integral descent assertion is required.

Finally, fix $a\in A_{S_D}$. Left multiplication by $a$ preserves the finite free $S_D$-module $M_{R_v}$. Its characteristic polynomial is monic over $S_D$, and Cayley--Hamilton, applied to the vector $1$, gives a monic polynomial over $S_D$ annihilating $a$. That polynomial also annihilates $L_a$ on $A_{F_D}$. Every eigenvalue of $L_a$ in an algebraic closure of $F_D$ is therefore integral over $S_D$. The coefficients of $\det(X-L_a)$ are elementary symmetric functions of these eigenvalues, so are integral over $S_D$. They belong to $F_D$, and $S_D=\mathbb Z[D]$ is integrally closed in $F_D$. Hence they belong to $S_D$, as claimed.
\end{proof}

The same calculation gives a finite certificate for a localization on which the specified word algebra is free.

\begin{corollary}\label{cor:hecke-word-localization}
Choose an $F_D$-word basis $w_1,\ldots,w_d$ of $A_{F_D}$. Express its vectors in the ambient $S_D$-basis
\[
 \bigl\{v^a(T_x\otimes T_y):0\le a<4,\quad x,y\in\mathfrak S_r\bigr\}.
\]
Let $f(D)\ne0$ be any $d$-by-$d$ minor of their coordinate matrix. Then
\[
 A_{S_D}[1/f]=\bigoplus_{i=1}^d S_D[1/f]w_i.
\]
In particular all multiplication coefficients in this word basis lie in $S_D[1/f]$.
\end{corollary}

\begin{proof}
The columns are independent over $F_D$, so a nonzero minor exists. Every word has coordinates in $S_D$ in the displayed ambient basis and belongs to the $F_D$-span of the selected words. Restrict its coordinate equations to the rows defining $f$. Cramer's rule puts each coefficient in $f^{-1}S_D\subseteq S_D[1/f]$. Conversely, each $w_i$ already lies in $A_{S_D}$. Localizing gives the asserted equality, and independence follows from the original field basis.
\end{proof}

At $r=4$, the theorem specifies a finite torsion-free integral word algebra and proves that its usual Laurent word algebra is its faithful scalar extension. The corollary supplies a denominator certificate once a word basis and a nonzero minor have been computed. These general arguments alone do not compute the minor, prove integral freeness, or construct a positive distinguished basis. The requirements are separate from the rank-four canonical-basis data: \citet[\S5.2]{Mulmuley} computes only the trivial-block basis element there. In particular this theorem does not identify a least coefficient ring for a positive rank-four basis, and it makes no assertion about descent of individual irreducible characters beyond the regular characteristic polynomials stated above.
